\documentclass[11pt,a4paper]{article}

\usepackage{amsmath}
\usepackage{amssymb}
\usepackage{geometry}
\usepackage{pgfplots}
\usepackage{pgfplotstable}
\usepackage{siunitx}
\usepackage{booktabs}
\usepackage{hyperref}
\usepackage{microtype}
\usepackage{parskip}

\pgfplotsset{compat=1.18}
\geometry{margin=2.5cm}

\sisetup{
  separate-uncertainty,
  per-mode=symbol
}

\hypersetup{
  colorlinks=true,
  linkcolor=black,
  urlcolor=blue
}

\title{Capacitance of a Tilted Circular Parallel-Plate Capacitor\\
       with Application to Precision Tilt Measurement}
\author{}
\date{}

\begin{document}

\maketitle
\tableofcontents
\bigskip

% ─────────────────────────────────────────────────────────────
\section{Problem Statement}
% ─────────────────────────────────────────────────────────────

Precision tilt sensors based on capacitive measurement exploit the
sensitivity of capacitance to small changes in the air gap between two
conducting plates.  In a practical instrument, a pendulum of length~$L$
suspends one plate so that its centre-to-centre distance to a fixed
reference plate varies with the tilt angle~$\alpha$.  When the
instrument is tilted, the two plates are no longer parallel: they subtend
the same angle~$\alpha$.

The goal of this note is to derive a closed-form expression for the
capacitance as a function of the tilt angle, compare it with two
progressively cruder approximations, and invert each expression to
obtain~$\alpha$ as a function of the measured capacitance.

% ─────────────────────────────────────────────────────────────
\section{Variable Definitions and Conventions}
% ─────────────────────────────────────────────────────────────

\begin{center}
\begin{tabular}{@{}lll@{}}
\toprule
Symbol & Unit & Description \\
\midrule
$\alpha$ & \si{\radian} & Tilt angle; $\alpha = 0$ corresponds to parallel plates \\
$L$      & \si{\metre}  & Pendulum length (constant) \\
$r$      & \si{\metre}  & Radius of the circular capacitor plates (constant) \\
$g$      & \si{\metre}  & Neutral air gap at $\alpha = 0$ (constant, $g > 0$) \\
$d_c$    & \si{\metre}  & Centre-to-centre plate distance at angle $\alpha$ \\
$\varepsilon_0$ & \si{\farad\per\metre} & Permittivity of free space,
                  $\varepsilon_0 = \SI{8.854e-12}{\farad\per\metre}$ \\
$C$      & \si{\farad}  & Total capacitance \\
$x$      & \si{\metre}  & Signed distance from the plate centre along the tilt axis \\
\bottomrule
\end{tabular}
\end{center}

By convention, $\alpha > 0$ means the pendulum swings in the positive
$x$-direction, closing the gap on that side.  Negative values of~$\alpha$
are equally valid and correspond to the opposite direction.

% ─────────────────────────────────────────────────────────────
\section{Geometry and Physical Limits}
% ─────────────────────────────────────────────────────────────

\subsection{Centre distance as a function of tilt}

The pendulum geometry gives the centre distance

\begin{equation}
  d_c(\alpha) = g - L\sin\alpha .
  \label{eq:dc}
\end{equation}

At $\alpha = 0$ the plates are parallel with uniform gap $g$.  For
$\alpha \neq 0$ the plates are tilted by the same angle~$\alpha$, so the
local gap at a signed distance~$x$ from the plate centre (measured along
the tilt axis) is

\begin{equation}
  d(x,\alpha) = d_c(\alpha) + x\sin\alpha .
  \label{eq:localGap}
\end{equation}

\subsection{Physical validity limit}

The plates must not touch anywhere.  The smallest gap occurs at
$x = -r$ (the trailing edge when $\alpha > 0$), so the condition
$d(-r,\alpha) > 0$ gives

\begin{equation}
  \boxed{\alpha < \arcsin\!\left(\frac{g}{L + r}\right).}
  \label{eq:alphaMax}
\end{equation}



% ─────────────────────────────────────────────────────────────
\section{Exact Capacitance}
% ─────────────────────────────────────────────────────────────

\subsection{From the parallel-plate formula to an integral}

The standard parallel-plate capacitor formula,

\begin{equation}
  C = \varepsilon_0 \frac{A}{d},
  \label{eq:parallelPlate}
\end{equation}

assumes a uniform gap~$d$ over the entire plate area~$A$.  When the gap
varies with position, equation~\eqref{eq:parallelPlate} is applied
locally to each infinitesimal area element $\mathrm{d}A$, and the total
capacitance is obtained by integrating the local contribution
$\mathrm{d}C = \varepsilon_0\,\mathrm{d}A / d(x,y)$:

\begin{equation}
  C = \varepsilon_0 \iint_{\text{disk}} \frac{\mathrm{d}A}{d(x,y)}.
  \label{eq:doubleIntegral}
\end{equation}

A double integral is required because the integration domain is a
two-dimensional disk.  Since the gap~\eqref{eq:localGap} depends only
on~$x$ and not on~$y$, the inner integral over~$y$ can be evaluated
immediately.  For the circular plate of radius~$r$ the $y$-limits at
a given~$x$ are $\pm\sqrt{r^2 - x^2}$, giving

\begin{equation}
  C(\alpha) = 2\varepsilon_0 \int_{-r}^{+r}
              \frac{\sqrt{r^2 - x^2}}{d_c + x\sin\alpha}\,\mathrm{d}x .
  \label{eq:singleIntegral}
\end{equation}

\subsection{Evaluation of the integral}

Equation~\eqref{eq:singleIntegral} is of the standard form

\begin{equation}
  I = \int_{-r}^{+r} \frac{\sqrt{r^2 - x^2}}{a + bx}\,\mathrm{d}x
    = \frac{\pi}{b^2}\!\left(a - \sqrt{a^2 - b^2 r^2}\right),
    \qquad a > br,
  \label{eq:stdIntegral}
\end{equation}

with $a = d_c$ and $b = \sin\alpha$.  Substituting and simplifying gives

\begin{equation}
  C(\alpha) = \frac{2\pi\varepsilon_0}{(\sin\alpha)^2}
              \!\left(d_c - \sqrt{d_c^2 - r^2(\sin\alpha)^2}\right).
  \label{eq:exactRaw}
\end{equation}

This form is indeterminate ($0/0$) at $\alpha = 0$.  Rationalising by
multiplying numerator and denominator by
$d_c + \sqrt{d_c^2 - r^2(\sin\alpha)^2}$ eliminates the singularity:

\begin{equation}
  \boxed{C(\alpha) = \frac{2\pi\varepsilon_0\, r^2}
         {d_c + \sqrt{d_c^2 - r^2(\sin\alpha)^2}},}
  \label{eq:exact}
\end{equation}

with $d_c = g - L\sin\alpha$ from equation~\eqref{eq:dc}.
Equation~\eqref{eq:exact} is valid for all~$\alpha$ satisfying the
physical limit~\eqref{eq:alphaMax}.

\subsection{Verification at \texorpdfstring{$\alpha = 0$}{alpha = 0}}

At $\alpha = 0$, $d_c \to g$ and $\sin\alpha \to 0$.  Expanding the
square root for small $s = \sin\alpha$,

\begin{equation}
  \sqrt{d_c^2 - r^2 s^2} \approx d_c - \frac{r^2 s^2}{2 d_c},
\end{equation}

so the denominator of~\eqref{eq:exact} approaches $2d_c \to 2g$, giving

\begin{equation}
  C(0) = \frac{\pi r^2 \varepsilon_0}{g} = \varepsilon_0 \frac{A}{g},
  \label{eq:C0}
\end{equation}

which is the standard parallel-plate result with $A = \pi r^2$. \checkmark

% ─────────────────────────────────────────────────────────────
\section{Approximate Formulae}
% ─────────────────────────────────────────────────────────────

\subsection{Approximation neglecting the tilt angle}

If the plate radius~$r$ is small compared with the ratio $d_c/\sin\alpha$,
the term $r^2(\sin\alpha)^2$ under the square root in~\eqref{eq:exact} is
negligible.  Physically, this corresponds to treating the plates as
parallel (using only the centre gap) while still accounting for the
change in $d_c$ with~$\alpha$.  Setting $r\sin\alpha \approx 0$ in the
denominator of~\eqref{eq:exact} reduces it to $2d_c$:

\begin{equation}
  \boxed{C_{\text{approx}}(\alpha) = \frac{\pi r^2 \varepsilon_0}{d_c}
         = \frac{\pi r^2 \varepsilon_0}{g - L\sin\alpha}.}
  \label{eq:approx}
\end{equation}

\subsection{Linear approximation}

A further simplification is obtained by linearising around $\alpha = 0$.
The sensitivity is found by differentiating~\eqref{eq:approx} with respect
to~$C$ at the neutral point.  Both the exact and approximate formulae
coincide at $\alpha = 0$, so their sensitivities are identical there.
Implicit differentiation gives

\begin{equation}
  \frac{\mathrm{d}\alpha}{\mathrm{d}C}\bigg|_{\alpha=0}
  = \frac{g^2}{\pi r^2 \varepsilon_0\, L},
\end{equation}

and the linear approximation around $\alpha = 0$, $C = C_0$ is

\begin{equation}
  \boxed{\alpha_{\text{lin}} = \frac{g^2}{\pi r^2 \varepsilon_0\, L}
         \!\left(C - C_0\right),}
  \label{eq:linear}
\end{equation}

with $C_0 = \pi r^2 \varepsilon_0 / g$ from equation~\eqref{eq:C0}.

% ─────────────────────────────────────────────────────────────
\section{Inversion: Solving for \texorpdfstring{$\alpha$}{alpha}}
% ─────────────────────────────────────────────────────────────

In practice, the capacitance~$C$ is the measured quantity and the tilt
angle~$\alpha$ is the desired output.  We solve each formula for~$\alpha$.

\subsection{Linear approximation}

Equation~\eqref{eq:linear} is already solved for~$\alpha$:

\begin{equation}
  \alpha_{\text{lin}} = \frac{g^2}{\pi r^2 \varepsilon_0\, L}
                        \!\left(C - C_0\right).
  \label{eq:alphaLin}
\end{equation}

\subsection{Approximate formula}

Solving~\eqref{eq:approx} for $\sin\alpha$ and applying the inverse sine:

\begin{equation}
  \boxed{\alpha_{\text{approx}} = \arcsin\!\left[\frac{1}{L}
         \!\left(g - \frac{\pi r^2 \varepsilon_0}{C}\right)\right].}
  \label{eq:alphaApprox}
\end{equation}

\subsection{Exact formula}

Starting from~\eqref{eq:exact} and writing $s = \sin\alpha$,
$d_c = g - Ls$, and $k = 2\pi\varepsilon_0 r^2 / C$:

\begin{align}
  d_c + \sqrt{d_c^2 - r^2 s^2} &= k \notag \\
  \sqrt{d_c^2 - r^2 s^2}        &= k - d_c \notag \\
  d_c^2 - r^2 s^2               &= k^2 - 2k\,d_c + d_c^2.
\end{align}

The $d_c^2$ terms cancel; substituting $d_c = g - Ls$ and collecting
terms in~$s$ yields a quadratic:

\begin{equation}
  r^2 s^2 + 2kL\,s + (k^2 - 2kg) = 0,
\end{equation}

with solution

\begin{equation}
  s = \frac{-kL \pm \sqrt{k^2(L^2 - r^2) + 2kgr^2}}{r^2},
  \qquad k = \frac{2\pi\varepsilon_0 r^2}{C}.
  \label{eq:quadSol}
\end{equation}

The two roots correspond to positive and negative tilt directions.
Only the root satisfying the physical limit~\eqref{eq:alphaMax} is
physically meaningful, so

\begin{equation}
  \boxed{\alpha_{\text{exact}} = \arcsin\!\left(
         \frac{-kL \pm \sqrt{k^2(L^2 - r^2) + 2kgr^2}}{r^2}
         \right),}
  \label{eq:alphaExact}
\end{equation}

where the sign is chosen such that $|\alpha| <
\arcsin\!\bigl(g/(L+r)\bigr)$.

% ─────────────────────────────────────────────────────────────
\section{Numerical Example and Graphical Comparison}
% ─────────────────────────────────────────────────────────────

\subsection{Instrument parameters}

\begin{center}
\begin{tabular}{@{}lll@{}}
\toprule
Parameter & Value & Description \\
\midrule
$L$ & \SI{40}{\milli\metre} & Pendulum length \\
$r$ & \SI{10}{\milli\metre} & Plate radius \\
$g$ & \SI{0.1}{\milli\metre} & Neutral air gap \\
\bottomrule
\end{tabular}
\end{center}

From equation~\eqref{eq:C0} the capacitance at $\alpha = 0$ is

\begin{equation}
  C_0 = \frac{\pi \times (10^{-2})^2 \times 8.854\times10^{-12}}{10^{-4}}
      = \SI{27.82}{\pico\farad}.
\end{equation}

The physical validity limit from equation~\eqref{eq:alphaMax} is

\begin{equation}
  \alpha_{\max} = \arcsin\!\left(\frac{0.1}{50}\right)
                = \arcsin(0.002) \approx \SI{412.5}{\arcsecond}
                \approx \SI{2000}{\micro\metre\per\metre}.
\end{equation}

\subsection{Comparison of the three models}

Figure~\ref{fig:comparison} shows $C$ as a function of~$\alpha$ for the
exact formula~\eqref{eq:exact}, the tilt-neglecting
approximation~\eqref{eq:approx}, and the linear
approximation~\eqref{eq:linear} over 70\,\% of the valid range.  A
secondary axis expresses the angle in \si{\micro\metre\per\metre}
(slope units used in metrology), related to arc seconds by
$1'' \approx 4.848\;\si{\micro\metre\per\metre}$.

\begin{figure}[ht]
\centering
\pgfplotstableread{capdata.dat}{\capdata}
\begin{tikzpicture}
\begin{axis}[
  width  = 0.85\textwidth,
  height = 7.5cm,
  xlabel = {Tilt angle $\alpha$ (arc seconds)},
  ylabel = {Capacitance $C$ (pF)},
  xmin = -300, xmax = 300,
  xtick distance = 100,
  grid = both,
  grid style = {line width=0.3pt, draw=gray!30},
  legend pos = north west,
  legend style = {font=\small},
  tick label style = {font=\small},
  label style = {font=\small},
  % Secondary x-axis (µm/m = arcsec * 4.8481)
  extra x ticks = {-288.77, -192.51, -96.26, 0, 96.26, 192.51, 288.77},
  extra x tick labels = {$-1400$, $-933$, $-467$, $0$, $467$, $933$, $1400$},
  extra x tick style = {
    xticklabel style = {yshift=0.65cm, font=\small},
    tick style = {yshift=0.4cm},
    grid = none,
  },
]
\addplot[blue,   thick,              ] table[x=arcsec, y=Cexact  ]{\capdata};
\addplot[teal,   thick, dashed,      ] table[x=arcsec, y=Capprox ]{\capdata};
\addplot[orange, thick, dotted, very thick] table[x=arcsec, y=Clinear ]{\capdata};
\legend{Exact \eqref{eq:exact},
        Approximate \eqref{eq:approx},
        Linear \eqref{eq:linear}}
\end{axis}
% µm/m label above the plot
\node at (7.2, 7.85) [font=\small] {Tilt (\si{\micro\metre\per\metre})};
\end{tikzpicture}
\caption{Capacitance versus tilt angle for the three models.
         Instrument parameters: $L=\SI{40}{\milli\metre}$,
         $r=\SI{10}{\milli\metre}$, $g=\SI{0.1}{\milli\metre}$.
         The plot covers 70\,\% of the physical validity range.
         The exact and approximate curves are nearly indistinguishable
         at this scale, confirming that the tilt correction is negligible
         for these dimensions.  The linear model diverges noticeably
         toward the edges of the range.}
\label{fig:comparison}
\end{figure}

% ─────────────────────────────────────────────────────────────
\section{Summary}
% ─────────────────────────────────────────────────────────────

Starting from the parallel-plate capacitance formula, three expressions
for the capacitance of a pendulum-coupled circular capacitor have been
derived and compared.

\begin{enumerate}

  \item \textbf{Exact formula} (eq.~\ref{eq:exact}): valid for all
        angles within the physical limit~\eqref{eq:alphaMax}; accounts
        fully for both the change in centre gap and the angular tilt of
        the plates.

  \item \textbf{Approximate formula} (eq.~\ref{eq:approx}): treats the
        plates as parallel using only the centre gap; accurate when
        $r\sin\alpha \ll d_c$.  For the example parameters the error is
        negligible across the entire operating range.

  \item \textbf{Linear approximation} (eq.~\ref{eq:linear}): a
        first-order Taylor expansion around $\alpha = 0$; valid only for
        small deviations from the neutral position.

\end{enumerate}

Each formula has been inverted analytically to yield $\alpha$ as a
function of~$C$ (equations~\ref{eq:alphaExact},
\ref{eq:alphaApprox}, \ref{eq:alphaLin}).  The exact inversion requires
solving a quadratic in $\sin\alpha$; both roots are real and correspond
to tilt in opposite directions, and the physically valid root is selected
by the limit~\eqref{eq:alphaMax}.


\end{document}
